\chapter{Testing and Evaluation}
This chapter presents how \sysname{} was tested and how well its AI components perform. Testing
checks that the software does what the requirements say (unit, integration, system, performance and
security testing). Evaluation measures the quality of the AI components on public, expert-labelled
datasets, because an AI system can pass every functional test and still give wrong answers.

\section{Testing Strategy}
\begin{itemize}
\item \textbf{Unit and integration tests} are automated with Django's test framework and run with
\code{python manage.py test} against a temporary PostgreSQL database with pgvector. They run without any
LLM (the LLM is switched off in the test settings), so they are fast, free and repeatable; the language
model is evaluated separately.
\item \textbf{System (black-box) tests} exercise complete use cases through the user interface or its REST API.
\item \textbf{Performance tests} measure response and processing times on the development machine.
\item \textbf{Security tests} check authentication, the role $\times$ access-level matrix, file access and
the audit log.
\item \textbf{AI evaluation} uses CUAD~\cite{hendrycks2021cuad}, LEDGAR~\cite{tuggener2020ledgar}, a
labelled set of 100 routing questions and a handwriting set (Section~\ref{sec:eval}).
\end{itemize}

\section{Unit Testing}
All automated tests pass (34 tests, 97 seconds). Table~\ref{tab:unit} lists the unit tests.

\begin{longtable}{|p{1.15cm}|p{3.1cm}|p{4.1cm}|p{4.5cm}|p{1.2cm}|}
\caption{Unit Test Cases}\label{tab:unit}\\ \hline
\textbf{ID} & \textbf{Module / Function} & \textbf{Input} & \textbf{Expected Output} & \textbf{Result} \\ \hline
\endfirsthead
\hline
\textbf{ID} & \textbf{Module / Function} & \textbf{Input} & \textbf{Expected Output} & \textbf{Result} \\ \hline
\endhead
UT-01 & \code{parse\_date} & ``2025-03-12'', ``12th day of March, 2025'', ``March 12, 2025'', ``12/03/2025'', ``sometime next year'' & First four parse to 12 March 2025 (day-first); the last returns none & Pass \\ \hline
UT-02 & \code{ground} & Page ``The parties signed this agreement on 12 March 2025 in Lahore.''; a real and an invented quote & Real quote grounded; ``signed on 1 January 2020 in Karachi'' rejected & Pass \\ \hline
UT-03 & \code{chunk\_pages} & Two Contract Act pages with ``CHAPTER II'' and numbered sections 10 and 11 & Section 11 chunk starts on page 2 and its path contains ``CHAPTER II'' and ``11. Who are competent'' & Pass \\ \hline
UT-04 & \code{visible\_documents} & Every role $\times$ ORG / DEPARTMENT / RESTRICTED documents & Each user sees exactly the allowed set & Pass \\ \hline
UT-05 & Documents API & Employee of another department requests a department-only document & Not listed; detail and file return 404 & Pass \\ \hline
UT-06 & Documents API & Admin deletes a document & Document is archived (row kept) and disappears from other users' lists & Pass \\ \hline
UT-07 & \code{can\_edit} & Employee edits someone else's document & HTTP 403 & Pass \\ \hline
UT-08 & Upload permission & Auditor uploads a file & HTTP 403 & Pass \\ \hline
UT-09 & Upload validation & \code{virus.exe}; a \code{.pdf} file whose content is not a PDF & Both rejected with HTTP 400 & Pass \\ \hline
UT-10 & Versions \& compare & Upload version 2, compare with version 1 & Two versions kept; added and removed lines returned & Pass \\ \hline
UT-11 & Field correction & Manager corrects a field value & Value and status change; FieldEdit with old and new value & Pass \\ \hline
UT-12 & Verify permission & Employee verifies a field & HTTP 403 & Pass \\ \hline
UT-13 & AuditEvent & Update and delete an existing event & Both raise an error; event unchanged & Pass \\ \hline
UT-14 & \code{heuristic\_plan} & ``When was the NDA with Acme signed?'' & LOOKUP, field signed\_date, party Acme & Pass \\ \hline
UT-15 & \code{\_clean} & ``null'', `` None '', ``Acme'' & None, None, ``Acme'' & Pass \\ \hline
UT-16 & \code{heuristic\_plan} & ``Find the lease for the Lahore office''; ``What makes an agreement void under the law?'' & SEARCH; KNOWLEDGE & Pass \\ \hline
UT-17 & \code{answer} (lookup) & ``When was the NDA with Acme signed?'' & Answer contains ``12 March 2025''; source is the NDA & Pass \\ \hline
UT-18 & Ask API \& history & Two questions in one conversation; another user opens it & Conversation has 2 queries; other user gets 404 & Pass \\ \hline
UT-19 & \code{notice\_days} & ``thirty (30) days written notice'', ``60 days'', ``one month'', ``2 weeks'', ``reasonable notice'' & 30, 60, 30, 14, none & Pass \\ \hline
UT-20 & \code{run\_reminders} & Lease expiring in 25 days with 30 days' notice & ``expires in 25 days'' and ``notice deadline 5 days ago'' reminders, including for the admin & Pass \\ \hline
UT-21 & \code{run\_reminders} & Contracts expiring in 400 and 60 days; check run twice & No reminder for the far one; exactly one per person for the near one & Pass \\ \hline
UT-22 & \code{people\_to\_remind} & HR department-only contract expiring in 5 days & Admin and uploader notified; legal employee and outsider not & Pass \\ \hline
UT-23 & Notifications API & List and mark as read & Unread count drops to 0 & Pass \\ \hline
UT-24 & Streaming ask API & ``When was the NDA with Acme signed?'' & Events arrive as plan, sources, text, done; the text contains ``12 March 2025''; the question is saved to history & Pass \\ \hline
UT-25 & Streaming ask API & HR employee asks for the secret settlement codeword, passing its document id & No source from the forbidden document and none of its text in the stream & Pass \\ \hline
UT-26 & Answer cache & Same question asked twice by one user, then by a user with different access & Second answer comes from the cache; the other user gets a fresh answer & Pass \\ \hline
UT-27 & \code{confident\_search} & ``Find the lease \dots'', ``Show me all handwritten affidavits'', ``Show me the expiry dates \dots'' & First two skip the LLM routing call; requests naming a record field do not & Pass \\ \hline
UT-28 & \code{parse\_term} & ``three (3) years'', ``thirty-six months'', ``90 days'', ``until terminated'' & 36 months, 36 months, 90 days, none & Pass \\ \hline
UT-29 & \code{derive\_expiry} & Effective 1 April 2025, term ``3 years'', no written expiry & Expiry 31 March 2028, quote = the term, status NEEDS\_REVIEW & Pass \\ \hline
UT-30 & \code{derive\_expiry} & Same, but with a written expiry date & Written date kept, nothing calculated & Pass \\ \hline
UT-31 & Second check & Governing law (supported) and a notice period quoting a renewal clause (not supported) & Governing law stays AUTO; notice period goes to review with the reason & Pass \\ \hline
\end{longtable}

\section{Integration Testing}
Integration tests check that modules work together. IT-01, IT-03 and IT-04 are automated; IT-02, IT-05 and
IT-06 were run against the running system.

\begin{longtable}{|p{1.15cm}|p{3.4cm}|p{4.0cm}|p{4.4cm}|p{1.3cm}|}
\caption{Integration Testing Summary}\label{tab:integration}\\ \hline
\textbf{ID} & \textbf{Modules Integrated} & \textbf{Input} & \textbf{Expected Result} & \textbf{Result} \\ \hline
\endfirsthead
\hline
\textbf{ID} & \textbf{Modules Integrated} & \textbf{Input} & \textbf{Expected Result} & \textbf{Result} \\ \hline
\endhead
IT-01 & Upload API $\rightarrow$ job queue $\rightarrow$ extraction $\rightarrow$ chunking $\rightarrow$ embeddings $\rightarrow$ PostgreSQL & A generated multi-page PDF & Version READY; pages, chunks with embeddings and search vectors stored & Pass \\ \hline
IT-02 & Vue Assistant $\rightarrow$ Ask API $\rightarrow$ router $\rightarrow$ retrieval $\rightarrow$ LLM $\rightarrow$ viewer & ``When is a contract voidable under the Contract Act?'' on the 65-page Contract Act 1872 & Answer citing the Contract Act with pages; source opens the page & Pass (13.6\,s) \\ \hline
IT-03 & Access control $\rightarrow$ hybrid retrieval & A secret document visible only to the legal department; searches by an HR employee, an outsider and a legal employee & HR employee and outsider never receive its chunks; legal employee gets it first & Pass \\ \hline
IT-04 & Ask API $\rightarrow$ engine with explicit document ids & HR employee passes the secret document's id & Id ignored; no sources from it & Pass \\ \hline
IT-05 & LLM gateway $\rightarrow$ fallback provider & Main provider's daily quota exhausted (HTTP 429, retry after 20 min) & Gateway switches to the fallback at once and answers; later calls skip the main provider & Pass \\ \hline
IT-06 & Docker Compose: database, backend, worker, nginx web & \code{docker compose up} on an empty database, then \code{seed\_demo} & Migrations run; all four containers healthy; six documents processed; sign-in, lookup, streamed answer and protected file download work through nginx & Pass \\ \hline
\end{longtable}

\section{System Testing}
Black-box tests follow the use cases of Chapter~3, signed in as the demo users of each role. BB-08 to
BB-11 were run in the user interface; the others were run through the same REST API the interface calls,
using a script that only sees requests and responses (no access to the code or database). During BB-10 a
defect was found: for long passages, every short word of the passage was highlighted all over the page. The
viewer was changed to locate the passage as a whole in the page text, and the test was repeated.

\begin{longtable}{|p{1.15cm}|p{2.0cm}|p{4.6cm}|p{5.0cm}|p{1.3cm}|}
\caption{System (Black Box) Test Cases}\label{tab:blackbox}\\ \hline
\textbf{ID} & \textbf{Use Case} & \textbf{Steps} & \textbf{Expected Result} & \textbf{Result} \\ \hline
\endfirsthead
\hline
\textbf{ID} & \textbf{Use Case} & \textbf{Steps} & \textbf{Expected Result} & \textbf{Result} \\ \hline
\endhead
BB-01 & UC-1.1 Register & Register a new organization with its admin; sign in & Organization created (201); signed-in user has the Admin role & Pass \\ \hline
BB-02 & UC-1.2 Login & Sign in with a wrong, then the correct password & 401, then tokens issued; LOGIN events in the audit log & Pass \\ \hline
BB-03 & UC-1.3 Members & Admin creates HR and Legal departments, adds an HR Employee, uploads an organization-wide and a Legal-only document & Member created (201); the HR employee sees only the organization-wide document & Pass \\ \hline
BB-04 & UC-2.1 Upload & Upload an image-only (scanned) PDF and a handwritten page marked handwritten & Both READY (46 s and 6 s); pages read by OCR and by handwriting recognition & Pass \\ \hline
BB-05 & UC-2.2 Versions & Upload a lease, then version 2 with a new rent, and compare & Two versions listed; comparison shows the new rent & Pass \\ \hline
BB-06 & UC-3.1 Register & Filter the records register by ``lease'' & Only leases listed, with expiry dates (31 Dec 2026, 30 Jun 2027) & Pass \\ \hline
BB-07 & UC-3.2 Review & Manager corrects the first field in the review queue with a reason; Employee tries to verify & Field leaves the queue; history stores the reason; Employee gets 403 & Pass \\ \hline
BB-08 & UC-4.1 Lookup & Ask ``Which contracts expire in 2027?'' & Office lease listed with expiry 30 June 2027 and evidence & Pass (1.5\,s) \\ \hline
BB-09 & UC-4.2 Law book & Ask ``What does the Contract Act say about agreements without consideration?'' & Section 25 rule and its exceptions with citations & Pass \\ \hline
BB-10 & UC-4.3 Source & Open the first source of a Contract Act answer & Page 22 opens with only the cited passage highlighted (Figure~\ref{fig:scr-source}) & Pass \\ \hline
BB-11 & UC-4.4 History & Reopen an earlier conversation & Questions, answers and sources shown & Pass \\ \hline
BB-12 & UC-5.1 Audit & Download a file; Auditor filters the audit log by DOWNLOAD; Employee opens the audit log & Only DOWNLOAD events listed; Employee gets 403 & Pass \\ \hline
\end{longtable}

\section{Performance Testing}
Measurements were taken on the development laptop (CPU only) with the free tiers of the LLM providers.
During the measurement the main provider's daily token quota had been used by the evaluation runs, so
generated answers came from the fallback provider (Gemini).

\begin{table}[H]\centering
\caption{Performance Metrics}\label{tab:perf}
\begin{tabularx}{\textwidth}{|p{4.4cm}|p{3.2cm}|X|}\hline
\textbf{Metric} & \textbf{Requirement} & \textbf{Observed Result} \\ \hline
Lookup answer time & $\leq$ 3 s (PER-1) & 1.5 s (``Which contracts expire in 2027?'') \\ \hline
Knowledge answer time & $\leq$ 35 s (PER-2) & 11.5--13.6 s with a warm server; 22.8 s for the first question after a restart (model loading) \\ \hline
Streamed knowledge answer & -- & Sources shown after 2.6 s, first words after 4.1 s, complete after 4.4 s (Groq, warm server) \\ \hline
Upload response & Immediate (PER-3) & Returns after storing the file; processing continues in the background \\ \hline
Digital page extraction & -- & About 10 ms per page (65-page Contract Act read in 1.2 s) \\ \hline
Handwritten page reading & -- & 6.8 s for one handwritten receipt (vision model) \\ \hline
Classification / field extraction & -- & 1.1--1.8 s / 1.4--2.0 s per short document; 38.6 s per long CUAD contract \\ \hline
Chunking and embedding & -- & 0.1--0.2 s for 1--3 page documents; 26 s for the 65-page Contract Act (190 chunks, CPU) \\ \hline
Automated test suite & -- & 34 tests in 97 s \\ \hline
\end{tabularx}
\end{table}

\noindent\textbf{Observation.} On free LLM tiers the main limit is the provider's token quota, not the
system: when the quota was exhausted, answers first took up to 45 seconds while the system waited for
retries. The gateway was then changed to switch provider immediately when a provider reports a long
wait, and to skip that provider until its quota resets, which produced the times above (IT-05).

\section{Security Testing}
\begin{table}[H]\centering
\caption{Security Verification Summary}\label{tab:security}
\begin{tabularx}{\textwidth}{|p{4.5cm}|X|p{2.1cm}|}\hline
\textbf{Check} & \textbf{Method} & \textbf{Result} \\ \hline
Unauthenticated access & API calls without a token return 401 & Pass \\ \hline
Role $\times$ access-level matrix & UT-04 checks every role against every access level & Pass \\ \hline
Retrieval leakage & IT-03 / IT-04: forbidden chunks never returned, even when their ids are sent & Pass \\ \hline
File access & UT-05: files of invisible documents return 404 from the permission-checked endpoint & Pass \\ \hline
Privilege escalation & UT-07, UT-08, UT-12: edit, upload and verify blocked for the wrong roles & Pass \\ \hline
Audit integrity & UT-13: events cannot be updated or deleted & Pass \\ \hline
Password storage & Django PBKDF2 salted hashes; plain passwords never stored & Pass \\ \hline
Rate limiting & DRF throttling: 20 questions/min, 600 requests/hour per user & Configured \\ \hline
Secrets & API keys only in \code{.env}, which is git-ignored & Pass \\ \hline
\end{tabularx}
\end{table}

\section{AI Evaluation}\label{sec:eval}
The AI components were evaluated with scripts in \code{evaluation/} on public datasets labelled by
experts. Each experiment ID (EV-$n$) is referenced in the traceability matrix.

\subsection{EV-1: Clause Classification (LEDGAR)}\label{sec:eval-clauses}
Test set: 10,000 provisions, 100 clause types. Macro-F1 averages over all types, so rare types count as
much as common ones.

\begin{table}[H]\centering
\caption{Clause Classification Results on LEDGAR}\label{tab:ev-clauses}
\begin{tabularx}{\textwidth}{|X|r|r|r|r|}\hline
\textbf{Model} & \textbf{Training provisions} & \textbf{Test provisions} & \textbf{Accuracy} & \textbf{Macro-F1} \\ \hline
TF-IDF + logistic regression (baseline) & 10,000 & 10,000 & 81.7\% & 72.3\% \\ \hline
TF-IDF + logistic regression (baseline) & 60,000 & 10,000 & 87.7\% & 82.5\% \\ \hline
Fine-tuned all-MiniLM-L6-v2 (ours, CPU, 3 epochs) & 10,000 & 10,000 & 71.2\% & 45.9\% \\ \hline
Fine-tuned all-MiniLM-L6-v2 (ours, GPU, 5 epochs) & 60,000 & 10,000 & 85.1\% & 75.1\% \\ \hline
\textbf{Fine-tuned MiniLM + TF-IDF, averaged (ours, used by the system)} & \textbf{60,000} & \textbf{10,000} & \textbf{88.5\%} & \textbf{82.7\%} \\ \hline
Zero-shot LLM (\code{gpt-oss-120b}, no training) & 0 & 198 & 65.2\% & 44.0\% \\ \hline
\end{tabularx}
\end{table}

The zero-shot LLM was tested on a random sample of 198 test provisions because of the free-tier rate limit; 2 calls that failed at the API were excluded. It needs no training data, but with 100 closely related labels it is less accurate than both trained models, which supports fine-tuning a small local model for this task: it is also free and fast at inference time.

Training on all 60,000 provisions on a free Colab T4 GPU (13.2 minutes) raised the fine-tuned model from
71.2\% to 85.1\% accuracy and from 45.9\% to 75.1\% macro-F1; only 4 of the 100 clause types are still never
predicted. The laptop run had been under-trained: with 10,000 examples many rare classes had almost no
training data.

On its own, the fine-tuned model is still below the TF-IDF baseline. LEDGAR clause types are signalled
largely by vocabulary (``governing law'', ``indemnify''), which suits a keyword model, and MiniLM is a very
small encoder (22 million parameters) chosen for fast CPU inference. The two models, however, make different
mistakes: TF-IDF relies on exact wording, the transformer on meaning. The system therefore averages their
predicted probabilities (\code{ml/train\_ensemble.py}). The weight of the transformer (0.25) was chosen on
the 10,000 validation provisions only, and the test set was used once. In that run the combination reached
\textbf{88.5\% accuracy and 82.7\% macro-F1}, better than TF-IDF (87.8\% / 82.1\%) and the fine-tuned model
(85.1\% / 75.0\%) evaluated on the same test set. This combination labels clauses in the application. All
trained models clearly beat the zero-shot LLM (65.2\%), at no API cost.

\subsection{EV-2: Retrieval (CUAD)}
For 60 CUAD contracts, 344 questions were asked about clauses that lawyers had marked. A result is a
hit when the returned passage contains the lawyers' answer span. Two chunking methods and three
retrieval methods were compared (an ablation study).

\begin{table}[H]\centering
\caption{Retrieval Results on CUAD (60 contracts, 344 questions)}\label{tab:ev-retrieval}
\begin{tabularx}{\textwidth}{|X|l|r|r|r|r|}\hline
\textbf{Chunking} & \textbf{Retrieval} & \textbf{Hit@1} & \textbf{Hit@3} & \textbf{Hit@5} & \textbf{MRR@10} \\ \hline
Fixed 1,000 characters (v1) & Semantic only & 49.7\% & 72.4\% & 82.8\% & 0.634 \\ \hline
Fixed 1,000 characters (v1) & Keyword only & 49.4\% & 73.8\% & 84.3\% & 0.637 \\ \hline
Fixed 1,000 characters (v1) & Hybrid RRF & 61.3\% & 82.3\% & 88.7\% & 0.733 \\ \hline
Section-aware (ours) & Semantic only & 53.5\% & 77.0\% & 84.6\% & 0.668 \\ \hline
Section-aware (ours) & Keyword only & 51.7\% & 80.2\% & 85.8\% & 0.665 \\ \hline
\textbf{Section-aware (ours)} & \textbf{Hybrid RRF} & \textbf{61.6\%} & \textbf{86.0\%} & \textbf{93.0\%} & \textbf{0.746} \\ \hline
\end{tabularx}
\end{table}
Compared with the design of the first version (fixed chunks, semantic search only), the final design
raises Hit@1 from 49.7\% to 61.6\% and Hit@5 from 82.8\% to 93.0\%. Both changes help on their own:
hybrid retrieval adds about 10 points of Hit@1, and section-aware chunking adds about 4 points of Hit@5.

\subsection{EV-3: Record Field Extraction (CUAD)}
Forty CUAD contracts were processed by the extraction step (model \code{gpt-oss-20b}, 38.6 s per
contract). A prediction is correct when its quoted evidence overlaps the lawyers' span or its value
matches (same date, parties or jurisdiction).

\begin{table}[H]\centering
\caption{Field Extraction Results on CUAD (40 contracts)}\label{tab:ev-extract}
\begin{tabular}{|l|r|r|r|}\hline
\textbf{Field} & \textbf{Precision} & \textbf{Recall} & \textbf{F1} \\ \hline
Parties & 100.0\% & 100.0\% & 100.0\% \\ \hline
Signing date & 88.0\% & 59.5\% & 71.0\% \\ \hline
Effective date & 77.1\% & 84.4\% & 80.6\% \\ \hline
Expiry date & 100.0\% & 35.7\% & 52.6\% \\ \hline
Auto-renewal & 90.0\% & 56.2\% & 69.2\% \\ \hline
Notice period & 26.9\% & 70.0\% & 38.9\% \\ \hline
Governing law & 90.0\% & 77.1\% & 83.1\% \\ \hline
\textbf{All fields (micro)} & \textbf{80.7\%} & \textbf{71.7\%} & \textbf{75.9\%} \\ \hline
\end{tabular}
\end{table}
\begin{itemize}
\item \textbf{Evidence is real.} 175 of 176 quotes (99.4\%) were found word-for-word on the cited page,
so every value can be checked against its source.
\item \textbf{Label differences.} CUAD's \emph{Notice Period To Terminate Renewal} covers only notice to
stop a renewal, while our notice period is any termination notice, so many correct values count as false
positives. CUAD's \emph{Expiration Date} is often a duration (``continues for five years''), while the
system stores an explicit date only when one is written, which lowers recall.
\item \textbf{Limitation of the first review rule.} Only 2 of the 34 wrong predictions were sent to review,
because most wrong values carry a real quote that answers a slightly different question. Confidence and
grounding alone do not catch these. Two improvements address this (below).
\end{itemize}

\paragraph*{Improvement 1: a second check.} After extraction, one more LLM call is asked a narrower question for
each value: does this quote, on its own, prove the value \emph{as this field}? Values that fail go to the review
queue with the model's reason (e.g.\ ``the quote is about a cure period, not termination notice''). Run on the
same 176 saved predictions (\code{evaluation/eval\_verifier.py}, model \code{gpt-oss-20b}):

\begin{table}[H]\centering
\caption{Second Check on Extracted Values (176 predictions)}\label{tab:ev-verify}
\begin{tabular}{|l|r|r|}\hline
\textbf{Review rule} & \textbf{Wrong values flagged} & \textbf{Right values flagged} \\ \hline
Confidence and grounding only & 2 of 34 (5.9\%) & 7 of 142 (4.9\%) \\ \hline
\textbf{With the second check} & \textbf{11 of 34 (32.4\%)} & 17 of 142 (12.0\%) \\ \hline
\end{tabular}
\end{table}
The check catches more than five times as many wrong values, mostly notice periods whose quote was really
about a cure period or payment terms, at the cost of reviewers also confirming 10 more correct values. Several
of those are signing dates quoted from ``IN WITNESS WHEREOF \dots{} as of the date above'', a quote that indeed
does not contain the date, so checking them is reasonable. (The saved quotes are cut to 200 characters, so
inside the pipeline the check sees slightly more text.)

\paragraph*{Improvement 2: calculated expiry dates.} Many contracts state a term (``for three (3) years from the
Effective Date'') instead of an end date. When no end date is written, the system now calculates it from the
effective or signing date and the term (start + term $-$ 1 day), quotes the term as evidence and always sends
the value to review. A written expiry date is never replaced. The 30 contracts without a written expiry date
were extracted again and the expiry field was scored with and without the calculation on the same run
(\code{evaluation/eval\_expiry.py}):

\begin{table}[H]\centering
\caption{Expiry Date with and without the Calculation (40 contracts)}\label{tab:ev-expiry}
\begin{tabular}{|l|r|r|r|}\hline
\textbf{Expiry date} & \textbf{Precision} & \textbf{Recall} & \textbf{F1} \\ \hline
Written dates only & 100.0\% & 35.7\% & 52.6\% \\ \hline
\textbf{With calculated dates} & 94.7\% & \textbf{64.3\%} & \textbf{76.6\%} \\ \hline
\end{tabular}
\end{table}
Nine dates were calculated and eight were right; the wrong one belongs to a contract the lawyers marked as
having no fixed expiry, and like every calculated date it was sent to review.

\subsection{EV-4: Query Routing}
A set of 100 questions (40 LOOKUP, 35 KNOWLEDGE, 25 SEARCH) labelled by hand was routed by the keyword
rules and by the LLM router. ``Right field'' is the share of lookup questions for which the right record
field (e.g.\ signing date) was chosen.
\begin{table}[H]\centering\footnotesize
\caption{Query Routing Results (100 questions)}\label{tab:ev-router}
\begin{tabular}{|l|r|r|r|r|r|r|}\hline
\textbf{Router} & \textbf{Accuracy} & \textbf{Lookup F1} & \textbf{Knowl.\ F1} & \textbf{Search F1} & \textbf{Right field} & \textbf{Time} \\ \hline
Keyword rules (fallback) & 77.0\% & 74.0\% & 73.2\% & 88.9\% & 60.0\% & $<$1 ms \\ \hline
LLM router (\code{gpt-oss-20b}) & 94.0\% & 96.4\% & 94.4\% & 88.9\% & 100.0\% & 1163 ms \\ \hline
\end{tabular}
\end{table}
The LLM router is correct on 94 of 100 questions against 77 for the
keyword rules, and it chooses the right field for every lookup (100.0\% against
60.0\%). Every LOOKUP question was routed correctly; most remaining errors are
search requests treated as lookups or knowledge questions (2 and 3 of 25), which still
return the right documents as sources. During the run, 11 of the 100 calls hit the free-tier rate limit and were
answered by the keyword fallback, so the LLM figures are a slight underestimate. A routing call takes about
1.2 s.

\subsection{EV-5: Handwriting Recognition}
Character and word error rates (lower is better) on 30 pages (synthetic handwriting-style fonts) for Tesseract, the vision
model, and the pipeline's rule (vision model only when Tesseract is unsure). The pages were generated from handwriting-style fonts, which
are easier than real handwriting; the script \code{evaluation/eval\_handwriting.py --folder} repeats the test on real
scans with transcripts.
\begin{table}[H]\centering
\caption{Handwriting Recognition Results}\label{tab:ev-hand}
\begin{tabular}{|l|r|r|r|}\hline
\textbf{Reader} & \textbf{CER} & \textbf{WER} & \textbf{Time per page} \\ \hline
Tesseract OCR & 5.0\% & 17.2\% & 1.0 s \\ \hline
Vision-language model (\code{gemini/gemini-2.5-flash}) & 0.1\% & 0.4\% & 34.8 s$^{\ast}$ \\ \hline
Pipeline (5/30 pages sent to the vision model) & 2.6\% & -- & -- \\ \hline
\end{tabular}
\end{table}
Tesseract alone already reads these typed-looking ``handwriting'' fonts well (5.0\% character errors), while the
vision model is almost perfect (0.1\%). With the pipeline's rule, only the 5 pages where Tesseract was unsure were
sent to the vision model, which halves the error rate (2.6\%) while keeping most pages fast and free. Six vision
calls were refused by the free-tier rate limit; those pages are left out of the vision figures.
$^{\ast}$The vision time includes waiting for the rate limit; a single uncontended call took about 6--7 s
(Table~\ref{tab:perf}).


\subsection{EV-6: Answer Quality}
CUAD questions were asked through the full answer engine for 10 contracts: about 20 whose answer is in the
contract and 10 whose answer is not. Measures: answer judged correct against the lawyers' annotation,
citation hit rate, and abstention (saying ``not in the document'') when the clause is absent. Correctness is
graded by a second LLM call that sees the lawyers' annotation (LLM-as-judge). Both runs used Gemini
Flash-Lite, the system's configured fallback model, for answering and grading: a run with \code{gpt-oss-20b}
hit the free daily token limit part-way (19 of 31 answers were ``model unavailable'' messages), so it was not
used.

\paragraph*{Improving ``not in the document'' answers.} In the first run the model sometimes presented a related
clause as if it answered the question. The answering instructions were changed so that, when the sources do
not contain the answer, the reply must begin with ``The documents do not say'' and then mention what related
information they do contain, without presenting it as the answer. The automatic check for such replies was
also widened (it had missed wording such as ``no specific duration is explicitly stated''); with the wider
check alone, the first run already scores 9 of 10.

\begin{table}[H]\centering
\caption{Answer Quality Before and After the New Instructions}\label{tab:ev-answers}
\begin{tabular}{|l|r|r|}\hline
\textbf{Measure} & \textbf{Before} & \textbf{After (used)} \\ \hline
Correctly says ``not in the document'' when absent & 80.0\% (8/10) & \textbf{100.0\% (10/10)} \\ \hline
At least one citation contains the lawyers' span & 95.2\% & \textbf{100.0\%} \\ \hline
Answers that cite their sources & 100.0\% & 100.0\% \\ \hline
Wrongly says ``not in the document'' when present & 4.8\% (1/21) & 10.0\% (2/20) \\ \hline
Answer judged correct & 100.0\% & 85.0\% (17/20) \\ \hline
\end{tabular}
\end{table}
The new instructions make the system more careful: it no longer answers any question whose clause is absent,
and every answer cites the lawyers' passage. The price is two over-cautious replies in borderline cases (a
``minimum percentage'' clause for a \emph{Minimum Commitment} question, and a termination notice for a
non-renewal notice question), where the reply nevertheless quotes and cites the relevant clause. For a legal
records system an occasional over-cautious reply is safer than an invented answer, so the new instructions
are used. (One question in the second run failed because the provider was overloaded; it is counted as an
error, not as a wrong answer. A third, softer wording was also tried and rejected: it answered 4 of 10
absent-clause questions with related clauses.) Because the same model graded its own answers, the
correctness figures are estimates and should be read together with the citation hit rate, which is measured
directly against the lawyers' spans.

\section{Testing Summary}
\begin{table}[H]\centering
\caption{Test Summary}\label{tab:testsummary}
\begin{tabular}{|l|r|r|r|}\hline
\textbf{Test type} & \textbf{Cases} & \textbf{Passed} & \textbf{Failed} \\ \hline
Unit (automated) & 31 & 31 & 0 \\ \hline
Integration (3 automated, 3 manual) & 6 & 6 & 0 \\ \hline
System (black box) & 12 & 12 & 0 \\ \hline
Security checks & 9 & 9 & 0 \\ \hline
\end{tabular}
\end{table}
All automated unit and integration tests pass, and the security checks confirm that access rules hold at
every layer, including retrieval. The evaluation shows that the redesigned retrieval clearly outperforms
the first version's design and that extracted values are backed by real evidence, while it also
identifies honest limitations (a small clause model that does not yet beat a strong keyword baseline, and the review rule's blind spot)
that are discussed in the next chapter.
